%********************************************************************
% Appendix
%*******************************************************
% If problems with the headers: get headings in appendix etc. right
%\markboth{\spacedlowsmallcaps{Appendix}}{\spacedlowsmallcaps{Appendix}}
% \chapter{Appendix Test}
% Lorem ipsum at nusquam appellantur his, ut eos erant homero
% concludaturque. Albucius appellantur deterruisset id eam, vivendum
% partiendo dissentiet ei ius. Vis melius facilisis ea, sea id convenire
% referrentur, takimata adolescens ex duo. Ei harum argumentum per. Eam
% vidit exerci appetere ad, ut vel zzril intellegam interpretaris.
% \graffito{More dummy text.}

%Errem omnium ea per, pro congue populo ornatus cu, ex qui dicant
%nemore melius. No pri diam iriure euismod. Graecis eleifend
%appellantur quo id. Id corpora inimicus nam, facer nonummy ne pro,
%kasd repudiandae ei mei. Mea menandri mediocrem dissentiet cu, ex
%nominati imperdiet nec, sea odio duis vocent ei. Tempor everti
%appareat cu ius, ridens audiam an qui, aliquid admodum conceptam ne
%qui. Vis ea melius nostrum, mel alienum euripidis eu.

% \section{Appendix Section Test}
% Test: \autoref{tab:moreexample} (This reference should have a 
% lowercase, small caps \spacedlowsmallcaps{A} if the option 
% \texttt{floatperchapter} is activated, just as in the table itself
%  $\rightarrow$ however, this does not work at the moment.)

% \begin{table}[h]
%     \myfloatalign
%   \begin{tabularx}{\textwidth}{Xll} \toprule
%     \tableheadline{labitur bonorum pri no} & \tableheadline{que vista}
%     & \tableheadline{human} \\ \midrule
%     fastidii ea ius & germano &  demonstratea \\
%     suscipit instructior & titulo & personas \\
%     %postulant quo & westeuropee & sanctificatec \\
%     \midrule
%     quaestio philosophia & facto & demonstrated \\
%     %autem vulputate ex & parola & romanic \\
%     %usu mucius iisque & studio & sanctificatef \\
%     \bottomrule
%   \end{tabularx}
%   \caption[Autem usu id]{Autem usu id.}
%   \label{tab:moreexample}
% \end{table}

%Nulla fastidii ea ius, exerci suscipit instructior te nam, in ullum
%postulant quo. Congue quaestio philosophia his at, sea odio autem
%vulputate ex. Cu usu mucius iisque voluptua. Sit maiorum propriae at,
%ea cum primis intellegat. Hinc cotidieque reprehendunt eu nec. Autem
%timeam deleniti usu id, in nec nibh altera.




% \section{Another Appendix Section Test}
% Equidem detraxit cu nam, vix eu delenit periculis. Eos ut vero
% constituto, no vidit propriae complectitur sea. Diceret nonummy in
% has, no qui eligendi recteque consetetur. Mel eu dictas suscipiantur,
% et sed placerat oporteat. At ipsum electram mei, ad aeque atomorum
% mea. There is also a useless Pascal listing below: \autoref{lst:useless}.

% \begin{lstlisting}[float=b,language=Pascal,frame=tb,caption={A floating example (\texttt{listings} manual)},label=lst:useless]
% for i:=maxint downto 0 do
% begin
% { do nothing }
% end;
% \end{lstlisting}


%Ei solet nemore consectetuer nam. Ad eam porro impetus, te choro omnes
%evertitur mel. Molestie conclusionemque vel at, no qui omittam
%expetenda efficiendi. Eu quo nobis offendit, verterem scriptorem ne
%vix.


%%% Local Variables:
%%% mode: latex
%%% TeX-master: "../ClassicThesis"
%%% End:


% ---------------------------------------------------------
% APPENDIX STARTS HERE
% ---------------------------------------------------------

\appendix
\chapter{Mathematical Foundations of Diffusion Models}
\label{app:diffusion_math}

\vspace{1cm}

This appendix provides the exhaustive mathematical proofs and derivations for the discrete-time diffusion framework outlined conceptually in \autoref{sec:bg_diffusion}. 

\section{The Forward Process and Reparameterization}
\label{app:diffusion_forward}

The single-step forward transition is defined as a Gaussian distribution centered around the scaled previous timestep:
\begin{equation}
    q(x_t | x_{t-1}) = \mathcal{N}(x_t; \sqrt{1 - \beta_t} x_{t-1}, \beta_t I).
\end{equation}

Using standard properties of Gaussian distributions, we can express a sample $x_t$ via the reparameterization trick as:
\begin{equation}
    x_t = \sqrt{1-\beta_t}x_{t-1} + \sqrt{\beta_t}\epsilon_{t-1},
\end{equation}
where $\epsilon_{t-1} \sim \mathcal{N}(0,I)$. 

By recursively applying the reparameterization trick and leveraging the property that the sum of two independent Gaussian variables $\mathcal{N}(0, \sigma_1^2 I)$ and $\mathcal{N}(0, \sigma_2^2 I)$ results in a new Gaussian $\mathcal{N}(0, (\sigma_1^2 + \sigma_2^2)I)$, we can merge the independent noise components at each step. Let $\alpha_t = 1 - \beta_t$ and $\overline{\alpha}_{t} = \prod_{s=1}^{t}\alpha_{s}$. The equation expands as follows:
\begin{align}
    x_t &= \sqrt{\alpha_t}x_{t-1} + \sqrt{1-\alpha_t}\epsilon_{t-1} \nonumber \\
        &= \sqrt{\alpha_t}(\sqrt{\alpha_{t-1}}x_{t-2} + \sqrt{1-\alpha_{t-1}}\epsilon_{t-2}) + \sqrt{1-\alpha_t}\epsilon_{t-1} \nonumber \\
        &= \sqrt{\alpha_t \alpha_{t-1}}x_{t-2} + \sqrt{\alpha_t(1-\alpha_{t-1}) + (1-\alpha_t)}\bar{\epsilon} \nonumber \\
        &= \dots \nonumber \\
        &= \sqrt{\overline{\alpha}_{t}}x_{0} + \sqrt{1-\overline{\alpha}_{t}}\epsilon.
\end{align}
This allows us to sample $x_t$ directly from $x_0$ without iterating through intermediate steps.

\section{Derivation of the Variational Lower Bound (VLB)}
\label{app:diffusion_vlb}

To optimize the neural network $p_\theta$, we must construct a tractable upper bound on the negative log-likelihood of the data. By introducing the forward process $q(x_{1:T}|x_0)$ and applying Jensen's inequality, we obtain the VLB:
\begin{align}
    -\log p_\theta(x_0) &= -\log \mathbb{E}_{q(x_{1:T}|x_0)} \left[ \frac{p_\theta(x_{0:T})}{q(x_{1:T}|x_0)} \right] \nonumber \\
    &\le \mathbb{E}_{q(x_{1:T}|x_0)} \left[ -\log \frac{p_\theta(x_{0:T})}{q(x_{1:T}|x_0)} \right] = \mathcal{L}_{VLB}.
\end{align}

We decompose this bound by expanding the joint distributions according to their Markov properties:
\begin{align}
    p_\theta(x_{0:T}) &= p(x_T) \prod_{t=1}^T p_\theta(x_{t-1}|x_t), \\
    q(x_{1:T}|x_0) &= \prod_{t=1}^T q(x_t|x_{t-1}).
\end{align}

Substituting these into the bound yields:
\begin{equation}
    \mathcal{L}_{VLB} = \mathbb{E}_q \left[ -\log p(x_T) - \sum_{t=1}^T \log p_\theta(x_{t-1}|x_t) + \sum_{t=1}^T \log q(x_t|x_{t-1}) \right].
\end{equation}

We resolve the intractable forward transition term $q(x_t|x_{t-1})$ by applying Bayes' rule conditioned on $x_0$:
\begin{equation}
    q(x_t|x_{t-1}) = \frac{q(x_{t-1}|x_t, x_0) q(x_t|x_0)}{q(x_{t-1}|x_0)}.
\end{equation}

Substituting this into the forward process summation and isolating the $t=1$ term gives:
\begin{equation}
    \sum_{t=1}^T \log q(x_t|x_{t-1}) = \log q(x_1|x_0) + \sum_{t=2}^T \log \left( \frac{q(x_{t-1}|x_t, x_0) q(x_t|x_0)}{q(x_{t-1}|x_0)} \right).
\end{equation}

Expanding the logarithm inside the sum produces a telescoping series:
\begin{equation}
    \sum_{t=2}^T \log q(x_{t-1}|x_t, x_0) + \sum_{t=2}^T \log q(x_t|x_0) - \sum_{t=2}^T \log q(x_{t-1}|x_0).
\end{equation}

The last two terms collapse to $\log q(x_T|x_0) - \log q(x_1|x_0)$. Recombining this with the isolated $t=1$ term yields a simplified sum:
\begin{equation}
    \sum_{t=1}^T \log q(x_t|x_{t-1}) = \sum_{t=2}^T \log q(x_{t-1}|x_t, x_0) + \log q(x_T|x_0).
\end{equation}

Substituting this back into the full $\mathcal{L}_{VLB}$ equation allows us to group the terms into mathematically elegant Kullback-Leibler (KL) divergences:
\begin{equation}
    \mathcal{L}_{VLB} = \mathbb{E}_q \left[ \log \frac{q(x_T|x_0)}{p(x_T)} + \sum_{t=2}^T \log \frac{q(x_{t-1}|x_t, x_0)}{p_\theta(x_{t-1}|x_t)} - \log p_\theta(x_0|x_1) \right].
\end{equation}

\begin{equation}
    \mathcal{L}_{VLB} = \mathbb{E}_q \left[ L_T + \sum_{t=2}^T L_{t-1} + L_0 \right],
\end{equation}
where $L_{t-1} = D_{KL}(q(x_{t-1} | x_t, x_0) \parallel p_\theta(x_{t-1} | x_t))$.

\section{Analytical Solution of the KL Divergence}
\label{app:diffusion_kl}

When conditioned on the original uncorrupted image $x_0$, the true reverse step $q(x_{t-1}|x_t, x_0)$ within $L_{t-1}$ becomes a tractable Gaussian distribution $\mathcal{N}(x_{t-1}; \tilde{\mu}_t(x_t, x_0), \tilde{\beta}_t I)$. The analytical mean and variance of this distribution are derived as:
\begin{align}
    \tilde{\mu}_t(x_t, x_0) &= \frac{\sqrt{\overline{\alpha}_{t-1}}\beta_t}{1-\overline{\alpha}_t}x_0 + \frac{\sqrt{\alpha_t}(1-\overline{\alpha}_{t-1})}{1-\overline{\alpha}_t}x_t, \\
    \tilde{\beta}_t &= \frac{1-\overline{\alpha}_{t-1}}{1-\overline{\alpha}_t}\beta_t.
\end{align}

Assuming our network $p_\theta$ is parameterized as a Gaussian with fixed variance $\Sigma_t = \sigma_t^2 I$ and a learned mean $\mu_\theta(x_t, t)$, the KL divergence mathematically reduces to the expected Mean Squared Error (MSE) between their means:
\begin{equation}
    L_{t-1} = \mathbb{E}_q \left[ \frac{1}{2\sigma_t^2} ||\tilde{\mu}_t(x_t, x_0) - \mu_\theta(x_t, t)||_2^2 \right] + C.
\end{equation}

By expressing $x_0$ algebraically in terms of $x_t$ and the noise $\epsilon$:
\begin{equation}
    x_0 = \frac{1}{\sqrt{\overline{\alpha}_t}}(x_t - \sqrt{1 - \overline{\alpha}_t}\epsilon),
\end{equation}
and substituting this into the target mean $\tilde{\mu}_t$, we see the target mean is entirely dictated by the injected noise. We parameterize the model's mean to mirror this relationship using a noise-prediction network $\epsilon_\theta(x_t, t)$:
\begin{equation}
    \mu_\theta(x_t, t) = \frac{1}{\sqrt{\alpha_t}} \left( x_t - \frac{\beta_t}{\sqrt{1-\overline{\alpha}_t}} \epsilon_\theta(x_t, t) \right).
\end{equation}

Substituting this parameterization back into $L_{t-1}$ simplifies the exact KL divergence into a weighted MSE loss:
\begin{equation}
    L_{t-1} - C = \mathbb{E}_{x_0, \epsilon} \left[ \frac{\beta_t^2}{2\sigma_t^2 \alpha_t (1-\overline{\alpha}_t)} ||\epsilon - \epsilon_\theta(\sqrt{\overline{\alpha}_t}x_0 + \sqrt{1-\overline{\alpha}_t}\epsilon, t)||_2^2 \right].
\end{equation}
As noted in \autoref{sec:bg_diffusion}, discarding the time-dependent weights yields the highly successful $\mathcal{L}_{simple}$ objective used in modern implementations.

\section{Generalized Sampling: From DDPM to DDIM}
\label{app:diffusion_sampling}

While \autoref{sec:bg_diffusion} conceptually introduces the difference between DDPM and DDIM samplers, we can formalize this distinction by examining the generalized reverse sampling step introduced by \textcite{song_denoising_2022}. 

For a generalized non-Markovian forward process that preserves the exact same marginal distribution $q(x_t|x_0)$ as DDPM, the reverse sampling step to generate $x_{t-1}$ from $x_t$ is defined as:

\begin{align}
x_{t-1} &= \nonumber\\ &\sqrt{\bar{\alpha}_{t-1}} \underbrace{\left( \frac{x_t - \sqrt{1-\bar{\alpha}_t} \epsilon_\theta(x_t, t)}{\sqrt{\bar{\alpha}_t}} \right)}_{\text{Predicted } x_0} \nonumber \\ &+ \underbrace{\sqrt{1 - \bar{\alpha}_{t-1} - \sigma_t^2} \cdot \epsilon_\theta(x_t, t)}_{\text{Direction pointing to } x_t} \nonumber \\ &+\underbrace{\sigma_t \omega}_{\text{Random noise}}
\end{align}

where $\omega \sim \mathcal{N}(0, I)$ is standard Gaussian noise, and $\sigma_t$ controls the stochastic variance injected into the process. 

The choice of $\sigma_t$ dictates the sampling paradigm:
\begin{itemize}
    \item \textbf{DDPM:} By setting $\sigma_t = \sqrt{(1 - \bar{\alpha}_{t-1}) / (1 - \bar{\alpha}_t)} \sqrt{1 - \alpha_t}$, the process matches the original Markovian trajectory of DDPM \parencite{ho_denoising_2020}.
    \item \textbf{DDIM:} By setting the stochastic variance to zero ($\sigma_t = 0$), the random noise term $\sigma_t z$ vanishes. The term inside the square root simplifies, leaving a purely deterministic update rule:
    
    $$x_{t-1} = \sqrt{\bar{\alpha}_{t-1}} \left( \frac{x_t - \sqrt{1-\bar{\alpha}_t} \epsilon_\theta(x_t, t)}{\sqrt{\bar{\alpha}_t}} \right) + \sqrt{1 - \bar{\alpha}_{t-1}} \cdot \epsilon_\theta(x_t, t)$$
\end{itemize}

Because this DDIM update equation contains no stochasticity, the reverse process becomes an Euler integration step of an Ordinary Differential Equation (ODE). This deterministic nature allows the model to safely evaluate the network on a sub-sequence of timesteps, accelerating generation without requiring the model to be retrained.

% \chapter{Internal UNET Encoder}

% \begin{lstlisting}[float=b,frame=tb,caption={best AURC(t) per output block sub-module},label=lst:useless]
% output_blocks.0.0 best AURC: 11.76 at t=5
% output_blocks.0.1 best AURC: 11.82 at t=5
% output_blocks.1.0 best AURC: 11.87 at t=5
% output_blocks.1.1 best AURC: 11.74 at t=5
% output_blocks.2.0 best AURC: 11.96 at t=5
% output_blocks.2.1 best AURC: 11.86 at t=5
% output_blocks.2.2 best AURC: 11.87 at t=5
% output_blocks.3.0 best AURC: 11.98 at t=5
% output_blocks.3.1 best AURC: 11.95 at t=5
% output_blocks.4.0 best AURC: 12.13 at t=5
% output_blocks.4.1 best AURC: 12.09 at t=5
% output_blocks.5.0 best AURC: 12.22 at t=5
% output_blocks.5.1 best AURC: 12.12 at t=5
% output_blocks.5.2 best AURC: 12.16 at t=5
% \end{lstlisting}

% \begin{table}[!htbp]
% \centering
% \caption{ADM UNet \texttt{output\_blocks} structure for ImageNet $128\times128$,
%          extracted from \texttt{model.named\_modules()}.
%          The \texttt{.2} sub-block in blocks 2, 5, 8, and 11 is the up/downsampling
%          ResBlock. Blocks 9--14 operate without attention.
%          We also assign indexes to sub-modules inside each block before their name: for instance \texttt{0.Res} in \texttt{output\_block.0} will be identified as \texttt{output\_blocks.0.0}, the first layer in the first block.}
% \label{tab:adm-unet-decoder}
% \begin{tabular}{llcc}
% \toprule
% \textbf{Block} $\downarrow$ & \textbf{Sub-modules} $\rightarrow$ & \textbf{Spatial res.} & \textbf{Channels}\\
% \midrule
% output\_blocks.0  & 0.Res + 1.Attention & $8\times8$     & 1024  \\
% output\_blocks.1  & 0.Res + 1.Attention & $8\times8$     & 1024  \\
% output\_blocks.2  & 0.Res + 1.Attention + Upsample   & $8\times8$     & 1024 \\
% \midrule
% output\_blocks.3  & 0.Res + 1.Attention & $16\times16$   & 512 \\
% output\_blocks.4  & 0.Res + 1.Attention & $16\times16$   & 512 \\
% output\_blocks.5  & 0.Res + 1.Attention + Upsample   & $16\times16$   & 512 \\
% \midrule
% % output\_blocks.6  & 0.Res + 1.Attention & $32\times32$   & 512 \\
% \dots & \dots & \dots & \dots\\
% % output\_blocks.7  & 0.Res + 1.Attention & $32\times32$   & 512 \\
% % output\_blocks.8  & 0.Res + 1.Attention + Upsample   & $32\times32$   & 512 \\
% % \midrule
% % output\_blocks.9  & 0.Res & $64\times64$   & 256  \\
% % output\_blocks.10 & 0.Res & $64\times64$   & 256  \\
% % output\_blocks.11 & 0.Res + Upsample                    & $64\times64$   & 256 \\
% \midrule
% output\_blocks.12 & 0.Res & $128\times128$ & 256  \\
% output\_blocks.13 & 0.Res & $128\times128$ & 256  \\
% output\_blocks.14 & 0.Res & $128\times128$ & 256  \\
% % \midrule
% % \multicolumn{5}{l}{\textit{Output head}} \\
% % out.0             & GroupNorm + SiLU                       & $128\times128$ & 256  \\
% % out.2             & Conv2d $3\times3 \to \hat{\epsilon}$   & $128\times128$ & 6   \\
% \bottomrule
% \end{tabular}
% \end{table}